\documentclass[../../main.tex]{subfiles}
\begin{document}
\section{习题}\label{section:ch07-习题}

\begin{exercise}\label{exercise:ch07-01}
证明对任意 Lie 群 $G$，乘法映射 $m\colon G\times G\to G$ 是光滑淹没。[提示：使用局部截面。]
\end{exercise}

\begin{exercise}\label{exercise:ch07-02}
设 $G$ 是 Lie 群。
\begin{enumerate}[(a)]
\item 记 $m\colon G\times G\to G$ 为乘法映射。利用\cref{proposition:ch03-14}将 $T_{(e,e)}(G\times G)$ 等同于 $T_eG\oplus T_eG$，证明微分
\begin{align*}
\mathrm{d}m_{(e,e)}\colon T_eG\oplus T_eG\to T_eG
\end{align*}
由 $\mathrm{d}m_{(e,e)}(X,Y)=X+Y$ 给出。[提示：分别计算 $\mathrm{d}m_{(e,e)}(X,0)$ 和 $\mathrm{d}m_{(e,e)}(0,Y)$。]
\item 记 $i\colon G\to G$ 为取逆映射。证明 $\mathrm{d}i_e\colon T_eG\to T_eG$ 由 $\mathrm{d}i_e(X)=-X$ 给出。
\end{enumerate}
（用于第 203、522 页。）
\end{exercise}

\begin{exercise}\label{exercise:ch07-03}
我们对 Lie 群的定义包含乘法映射和取逆映射都光滑的要求。证明取逆映射的光滑性是多余的：若 $G$ 是具有群结构的光滑流形，且乘法映射 $m\colon G\times G\to G$ 是光滑的，则 $G$ 是 Lie 群。[提示：证明由 $F(g,h)=(g,gh)$ 定义的映射 $F\colon G\times G\to G\times G$ 是双射局部微分同胚。]
\end{exercise}

\begin{exercise}\label{exercise:ch07-04}
记 $\det\colon\operatorname{GL}(n,\mathbb{R})\to\mathbb{R}$ 为行列式函数。利用\cref{corollary:ch03-25}按如下方式计算 $\det$ 的微分。
\begin{enumerate}[(a)]
\item 对任意 $A\in\mathcal{M}(n,\mathbb{R})$，证明
\begin{align*}
\left.\frac{\mathrm{d}}{\mathrm{d}t}\right|_{t=0}\det(I_n+tA)=\operatorname{tr}A,
\end{align*}
其中 $\operatorname{tr}(A_i^j)=\sum_i A_i^i$ 是 $A$ 的迹。[提示：定义方程 (B.3) 将 $\det(I_n+tA)$ 表示为 $t$ 的多项式。线性项是什么？]
\item 对 $X\in\operatorname{GL}(n,\mathbb{R})$ 和 $B\in T_X\operatorname{GL}(n,\mathbb{R})\cong\mathcal{M}(n,\mathbb{R})$，证明
\begin{align}
\mathrm{d}(\det)_X(B)=(\det X)\operatorname{tr}(X^{-1}B).\label{align:ch07-11}
\end{align}
[提示：$\det(X+tB)=\det(X)\det(I_n+tX^{-1}B)$。]
\end{enumerate}
（用于第 203 页。）
\end{exercise}

\begin{exercise}\label{exercise:ch07-05}
证明\cref{theorem:ch07-09}（万有覆叠群的唯一性）。
\end{exercise}

\begin{exercise}\label{exercise:ch07-06}
设 $G$ 是 Lie 群，$U$ 是单位元的任意邻域。证明存在单位元的邻域 $V$ 使得 $V\subseteq U$，且当 $g,h\in V$ 时 $gh^{-1}\in U$。（用于第 159、556 页。）
\end{exercise}

\begin{exercise}\label{exercise:ch07-07}
证明\cref{proposition:ch07-15}（Lie 群单位分支的性质）。
\end{exercise}

\begin{exercise}\label{exercise:ch07-08}
设连通拓扑群 $G$ 连续作用于离散空间 $K$。证明该作用是平凡的。（用于第 562 页。）
\end{exercise}

\begin{exercise}\label{exercise:ch07-09}
证明公式 $A\cdot[x]=[Ax]$ 定义了 $\operatorname{GL}(n+1,\mathbb{R})$ 在 $\mathbb{RP}^n$ 上的光滑、传递左作用。
\end{exercise}

\begin{exercise}\label{exercise:ch07-10}
对 $\operatorname{GL}(n+1,\mathbb{C})$ 和 $\mathbb{CP}^n$ 重复\cref{exercise:ch07-09}（见\cref{exercise:ch01-09}和\cref{exercise:ch04-05}）。
\end{exercise}

\begin{exercise}\label{exercise:ch07-11}
将 $S^{2n+1}$ 视为 $\mathbb{C}^{n+1}$ 中的单位球面，定义 $S^1$ 在 $S^{2n+1}$ 上的作用，称为\textbf{Hopf 作用}，即
\begin{align*}
z\cdot(w^1,\dots,w^{n+1})=(zw^1,\dots,zw^{n+1}).
\end{align*}
证明该作用是光滑的，其轨道是 $\mathbb{C}^{n+1}$ 中不交的单位圆，它们的并为 $S^{2n+1}$。（用于第 560 页。）
\end{exercise}

\begin{exercise}\label{exercise:ch07-12}
利用等变秩定理给出\cref{theorem:ch07-05}的另一种证明，方法是证明每个 Lie 群同态 $F\colon G\to H$ 关于 $G$ 和 $H$ 上适当的光滑 $G$-作用是等变的。
\end{exercise}

\begin{exercise}\label{exercise:ch07-13}
对每个 $n\geqslant1$，证明 $U(n)$ 是 $\operatorname{GL}(n,\mathbb{C})$ 的维数为 $n^2$ 的真嵌入 Lie 子群。（见\cref{example:ch07-29}。）
\end{exercise}

\begin{exercise}\label{exercise:ch07-14}
对每个 $n\geqslant1$，证明 $\operatorname{SU}(n)$ 是 $U(n)$ 的维数为 $n^2-1$ 的真嵌入 Lie 子群。（见\cref{example:ch07-30}。）
\end{exercise}

\begin{exercise}\label{exercise:ch07-15}
证明 $\operatorname{SO}(2)$、$U(1)$ 和 $S^1$ 作为 Lie 群都是同构的。
\end{exercise}

\begin{exercise}\label{exercise:ch07-16}
证明 $\operatorname{SU}(2)$ 与 $S^3$ 微分同胚。（用于第 179、563 页。）
\end{exercise}

\begin{exercise}\label{exercise:ch07-17}
确定下列 Lie 群中哪些是紧的：
\begin{align*}
\operatorname{GL}(n,\mathbb{R});\quad \operatorname{SL}(n,\mathbb{R});\quad \operatorname{GL}(n,\mathbb{C});\quad \operatorname{SL}(n,\mathbb{C});\quad U(n);\quad \operatorname{SU}(n).
\end{align*}
\end{exercise}

\begin{exercise}\label{exercise:ch07-18}
证明\cref{theorem:ch07-35}（半直积的刻画）。
\end{exercise}

\begin{exercise}\label{exercise:ch07-19}
设 $G$、$N$ 和 $H$ 是 Lie 群。证明 $G$ 同构于半直积 $N\rtimes H$ 当且仅当存在 Lie 群同态 $\varphi\colon G\to H$ 和 $\psi\colon H\to G$ 使得 $\varphi\circ\psi=\operatorname{Id}_H$ 且 $\operatorname{Ker}\varphi\cong N$。
\end{exercise}

\begin{exercise}\label{exercise:ch07-20}
证明下列 Lie 群同构于所示的半直积。[提示：使用\cref{exercise:ch07-19}。]
\begin{enumerate}[(a)]
\item $O(n)\cong\operatorname{SO}(n)\rtimes O(1)$。
\item $U(n)\cong\operatorname{SU}(n)\rtimes U(1)$。
\item $\operatorname{GL}(n,\mathbb{R})\cong\operatorname{SL}(n,\mathbb{R})\rtimes\mathbb{R}^*$。
\item $\operatorname{GL}(n,\mathbb{C})\cong\operatorname{SL}(n,\mathbb{C})\rtimes\mathbb{C}^*$。
\end{enumerate}
\end{exercise}

\begin{exercise}\label{exercise:ch07-21}
证明当 $n>1$ 时，\cref{exercise:ch07-20}中的群不同构于所示群的直积。[提示：群 $G$ 的\textbf{中心}是与 $G$ 中每个元素都交换的所有元素的集合。证明同构群有同构的中心。]
\end{exercise}

\begin{exercise}\label{exercise:ch07-22}
令 $H=\mathbb{C}\times\mathbb{C}$（视为实向量空间），定义双线性乘积 $H\times H\to H$ 为
\begin{align*}
(a,b)(c,d)=(ac-d\overline{b},\,\overline{a}d+cb),
\end{align*}
对 $a,b,c,d\in\mathbb{C}$。在这个乘积下，$H$ 是 $\mathbb{R}$ 上的四维代数，称为\textbf{四元数代数}。对每个 $p=(a,b)\in H$，定义 $p^*=(\overline{a},-b)$。使用由
\begin{align*}
1=(1,0),\quad \mathrm{i}=(\mathrm{i},0),\quad j=(0,1),\quad k=(0,-\mathrm{i})
\end{align*}
定义的 $H$ 的基 $(1,\mathrm{i},j,k)$ 是有用的。直接验证这个基满足
\begin{align*}
\mathrm{i}^2=j^2=k^2=-1,\qquad 1q=q1=q\quad\text{对所有 }q\in H;\\
\mathrm{i}j=-j\mathrm{i}=k,\quad jk=-kj=\mathrm{i},\quad k\mathrm{i}=-\mathrm{i}k=j;\\
1^*=1,\quad \mathrm{i}^*=-\mathrm{i},\quad j^*=-j,\quad k^*=-k.
\end{align*}
四元数 $p$ 称为\textbf{实的}，若 $p^*=p$；称为\textbf{虚的}，若 $p^*=-p$。实四元数可通过对应 $x\leftrightarrow x1$ 与实数等同。
\begin{enumerate}[(a)]
\item 证明四元数乘法是结合的但不是交换的。
\item 证明对所有 $p,q\in H$ 有 $(pq)^*=q^*p^*$。
\item 证明 $\langle p,q\rangle=\frac{1}{2}(p^*q+q^*p)$ 是 $H$ 上的内积，其相伴范数满足 $|pq|=|p||q|$。
\item 证明每个非零四元数有双边乘法逆，由 $p^{-1}=|p|^{-2}p^*$ 给出。
\item 证明非零四元数集合 $H^*$ 在四元数乘法下是 Lie 群。
\end{enumerate}
（用于第 200、200、562 页。）
\end{exercise}

\begin{exercise}\label{exercise:ch07-23}
设 $H^*$ 为非零四元数的 Lie 群（\cref{exercise:ch07-22}），$\Sigma\subseteq H^*$ 为单位四元数的集合。证明 $\Sigma$ 是 $H^*$ 的真嵌入 Lie 子群，同构于 $\operatorname{SU}(2)$。（用于第 200、562 页。）
\end{exercise}

\begin{exercise}\label{exercise:ch07-24}
证明\cref{example:ch07-36}(g)中描述的每个映射 $\tau_d^n\colon\operatorname{GL}(n,\mathbb{R})\to\operatorname{GL}(\mathcal{P}_d^n)$ 都是 $\operatorname{GL}(n,\mathbb{R})$ 的忠实表示。
\end{exercise}

\end{document}